\documentclass[10pt,a4paper]{article}
\usepackage[margin=0.72in]{geometry}
\usepackage{xcolor}
\usepackage{enumitem}
\usepackage{fancyhdr}
\usepackage{hyperref}
\definecolor{navy}{HTML}{17324D}
\definecolor{teal}{HTML}{0D7468}
\definecolor{coral}{HTML}{B95843}
\definecolor{paper}{HTML}{F4F7F6}
\hypersetup{colorlinks=true,linkcolor=teal,urlcolor=teal,pdftitle={Chakra Feasibility and Idea Counter Questions}}
\pagestyle{fancy}\fancyhf{}\lhead{\textcolor{navy}{CHAKRA}}\rhead{\textcolor{teal}{dot\_gitignore}}\cfoot{\thepage}\setlength{\headheight}{14pt}
\setlist[enumerate]{leftmargin=*,itemsep=8pt,topsep=5pt}
\newcommand{\qa}[2]{\item \textbf{\color{navy}Q: #1}\\[-1pt]{\color{coral}\textit{Answer:}} #2}
\newcommand{\note}[1]{\vspace{4pt}\noindent\colorbox{paper}{\parbox{\dimexpr\linewidth-2\fboxsep}{\small\textbf{How to answer:} #1}}\vspace{6pt}}
\begin{document}
\begin{center}
{\Huge\bfseries\color{navy}Chakra}\\[3pt]
{\Large Feasibility, Idea and Impact Counter Questions}\\[5pt]
{\normalsize Smart India Hackathon SIH26189 \textbar\ Team dot\_gitignore}
\end{center}
\note{These answers describe the working prototype honestly: fictional women-safety records, a frontend demo, explainable graph leads, and human review. Replace any demo number with a measured result before the final pitch.}

\section{Problem, users and need}
\begin{enumerate}
\qa{What exact problem does Chakra solve?}{Chakra helps authorized investigators discover possible links between women-safety cases that are difficult to see across separate records. It organizes evidence by time, source and case, then presents a reviewable lead instead of making an accusation.}
\qa{Why is this a real problem for the current workflow?}{Important clues are split across FIRs, statements, call events, transactions and registries. Manual comparison is slow, repeated identities are easy to miss, and an analyst may not be able to explain how a cross-case connection was found.}
\qa{Who is the primary user?}{The primary user is an authorized investigator or analyst working on a case. Supervisors and report reviewers are secondary users who check decisions, access scope and evidence provenance.}
\qa{Why focus on women-safety cases?}{Women-safety investigations can involve repeat patterns, shared contacts, unsafe locations and coordinated assistance. Chakra is designed around those investigative needs while keeping survivor identity protected and avoiding automated guilt labels.}
\qa{What is the one-line value proposition?}{Chakra turns fragmented, time-stamped case evidence into explainable cross-case leads that an authorized officer can verify, accept, reject or correct.}
\qa{What makes the idea different from a normal case-management system?}{A case-management system stores and searches records. Chakra adds a temporal relationship layer that makes links between cases visible, shows the evidence path, and separates observed facts from inferred suggestions.}
\qa{Why is a graph the right representation?}{A graph represents people, phones, accounts, places, events and cases as entities with relationships. It makes bridge entities, communities and repeated patterns easier to inspect than a list of documents.}
\qa{Why is time essential to the idea?}{The same relationship can mean different things before and after an incident. Every event and edge therefore carries a timestamp or time range so an investigator can ask whether a connection existed at the relevant moment.}
\qa{What does Chakra actually recommend?}{It recommends a case-link lead such as a shared masked phone, a transaction recipient or a repeated location pattern. The recommendation includes the relationship, confidence, source references and a route for human review.}
\qa{Does Chakra identify a criminal?}{No. It identifies an investigative lead, never guilt. The interface labels observed evidence and AI-inferred relationships separately and requires an officer to review the source before a lead can appear in an auditable report.}
\end{enumerate}

\section{Feasibility and delivery}
\begin{enumerate}[resume]
\qa{Can this be built within a hackathon?}{Yes, if the scope stays focused: import a small set of structured and scanned records, extract a limited entity vocabulary, build a temporal graph, score a few lead types, and demonstrate the review workflow. Production integrations and certification would follow after the prototype.}
\qa{What is the minimum viable demonstration?}{Upload three fictional case records, extract names and contact tokens, connect a CDR or payment event, show a bridge entity, open its evidence path, and let the investigator approve or reject the lead.}
\qa{What is real and what is simulated in the demo?}{The workflow, graph interaction, evidence presentation and decision states are implemented in the prototype. The records are fictional and the backend connectors, model serving and production access controls are represented as architecture to be integrated later.}
\qa{What data is required to start?}{A case identifier, source document, timestamp, entity or event fields, legal access scope and a cryptographic hash are the minimum provenance fields. Optional sources include CDR events, transaction records, vehicle registries and legally authorized OSINT.}
\qa{How will the system work with Indian languages?}{The ingestion plan supports OCR and multilingual NLP for Hindi, English and regional text. The first release should measure language-specific extraction quality and route low-confidence text to an officer rather than silently normalizing it.}
\qa{What if the department cannot share live data during judging?}{The demo uses a synthetic corpus that preserves the structure of the problem without exposing personal information. The same schema can later accept an approved data extract through an on-premise connector.}
\qa{How will success be measured?}{Measure extraction F1, entity-resolution precision and recall, case-link Precision@K and Recall@K, provenance coverage, false-positive rate, review time and officer override rate. A useful lead must be both accurate enough and explainable enough to act on.}
\qa{What is the expected benefit to an investigator?}{Chakra reduces time spent searching across records and gives the investigator a direct route from a graph relationship to the exact supporting source. It improves prioritization while leaving the investigative decision with the officer.}
\qa{How will you prove that the idea is not just a visual demo?}{The prototype uses deterministic fixtures, a defined graph schema, explicit evidence references, lead states and tests for state transitions. A next pilot would replay labeled cases and compare Chakra's ranked leads with expert annotations.}
\qa{What is the implementation roadmap after the hackathon?}{First stabilize ingestion, provenance and access controls; then pilot entity resolution and graph analytics on a controlled data set; then integrate approved department systems and conduct security, language and usability evaluations before wider rollout.}
\end{enumerate}

\section{Trust, adoption and impact}
\begin{enumerate}[resume]
\qa{Why would investigators trust an AI-assisted lead?}{Trust comes from showing the path, source document, timestamp, confidence and reasoning behind each lead. The officer can reject or correct it, and every decision is recorded rather than hidden behind a single score.}
\qa{How do you prevent automation bias?}{The interface uses neutral language such as ``suggested connection'', separates observed from inferred information, displays uncertainty, and requires a review action. Training should explicitly state that a lead is not a finding of guilt.}
\qa{What happens when the system is wrong?}{The investigator can reject, merge or correct a lead. That feedback is retained as an auditable decision and can later be used to improve models after quality review.}
\qa{How does Chakra protect survivors and witnesses?}{Names and sensitive fields can be masked by role, survivor identities can be restricted, exports are watermarked, and a supervisor approval is required for sensitive joins. The demo uses fictional records only.}
\qa{Could this system create bias against a community or language group?}{Yes, if source coverage or entity resolution is uneven. Chakra therefore needs subgroup and language error analysis, human review, calibrated thresholds and a policy against using graph prominence as a proxy for culpability.}
\qa{What is the ethical boundary of the product?}{It supports evidence organization and investigative prioritization. It must not produce an automated guilt score, unreviewed accusation label, predictive policing list or decision about a person's legal status.}
\qa{How does the idea align with public-sector procurement?}{The architecture is modular, deployable on premises, offline-friendly for imports and exports, and based on auditable open data structures. A department can adopt ingestion, graph or assistant components in stages.}
\qa{What is the cost argument?}{The largest value is analyst time saved and earlier discovery of a verifiable link. A pilot can begin with existing infrastructure and a narrow use case, then justify scale using measured review-time reduction and lead quality rather than speculative savings.}
\qa{Can small districts use it, or is it only for a central agency?}{Both are possible. A district can use a local case graph, while an approved supervisory deployment can compare carefully scoped entities across districts without exposing unrelated case data.}
\qa{How will you avoid mission creep?}{Define permitted data sources, user roles, retention periods and allowed questions before deployment. New joins or new OSINT sources should require documented legal and supervisory approval.}
\qa{What would make the judges believe this can be adopted?}{The prototype shows a complete path from evidence to lead to human decision, not only a graph visualization. The design also acknowledges integration, privacy, language, evaluation and governance work that must happen before production.}
\qa{What is the strongest limitation of the current proposal?}{The demo does not yet prove performance on real, noisy, multilingual departmental data. We present that as the next validation milestone and avoid claiming production accuracy from synthetic records.}
\end{enumerate}

\section{Pitch, differentiation and difficult follow-ups}
\begin{enumerate}[resume]
\qa{Why should this be selected over a search engine or dashboard?}{Search returns records that contain a term; Chakra connects entities and events across records, ranks a reviewable relationship and preserves the evidence path. It is designed for investigative reasoning rather than only retrieval.}
\qa{Why not let an LLM answer the investigator directly?}{An unconstrained LLM can invent links or lose provenance. Chakra restricts the assistant to graph-backed evidence, displays source references and keeps the graph and document path as the primary verification surface.}
\qa{What if two people share the same name or phone?}{The system treats a match as a candidate identity resolution, not proof. It combines aliases, time, location, case context and additional identifiers, keeps alternatives visible and sends ambiguous matches to review.}
\qa{What if the most important clue is missing from the data?}{Chakra cannot infer an absent fact reliably. It should show coverage limits, mark a lead as incomplete and allow an officer to add an authorized source rather than filling the gap with speculation.}
\qa{What is the best live demo story?}{Start with three fictional women-safety cases, show one shared masked contact and one payment relationship, open the source lines and timestamps, then approve one lead and reject another. The contrast demonstrates both usefulness and restraint.}
\qa{What is the success condition for the hackathon prototype?}{A judge should be able to follow every demonstrated lead from a case to a relationship to its evidence, understand the uncertainty, and complete a human review without needing a backend operator.}
\qa{What question should the team ask the department after the pitch?}{Which cross-case link types consume the most analyst time today, and what access, retention and review rules must be satisfied before a pilot can use them? That answer should shape the first real integration.}
\qa{If you had one more month, what would you build first?}{A labeled replay harness with real-world-like multilingual noise, a hardened provenance service and a small investigator usability pilot. Those three activities test whether the concept works beyond the polished demo.}
\end{enumerate}
\vfill
\begin{center}\small\color{teal}Chakra is an evidence-led decision-support prototype. Every lead remains subject to officer verification.\end{center}
\end{document}
